\subsection{Macaulay 环上的 Grothendieck 对偶}

我们继续假设 A 是维数为 \(d\) 的局部 Macaulay 环。设 \(\Omega\) 为一个对偶化 A-模。A-模 \(\mathrm{H}_{\mathrm{A}}^{d}(\Omega)\) 是一个 Matlis 模（ \(\S 8\) , 第 5 节，例 6）；按照 \(\S 8\) 的记号，对每个 A-模 M，令 \(\mathrm{D(M)} = \mathrm{Hom}_{\mathrm{A}}(\mathrm{M}, \mathrm{H}_{\mathrm{A}}^{d}(\Omega))\) 。

考虑同构 \(\tau^{d}(\Omega):\Omega\otimes\mathrm{H}_{\mathrm{A}}^{d}(\mathrm{A})\to\mathrm{H}_{\mathrm{A}}^{d}(\Omega)\) （第 2 节）。由此导出一个同态 \(\omega:\mathrm{H}_{\mathrm{A}}^{d}(\mathrm{A})\to\mathrm{Hom}_{\mathrm{A}}(\Omega,\mathrm{H}_{\mathrm{A}}^{d}(\Omega))\) ，它把每个元素 \(u\) （属于 \(\mathrm{H}_{\mathrm{A}}^{d}(\mathrm{A})\) ）对应到同态 \(x\mapsto\mathrm{H}_{\mathrm{A}}^{d}(f_{x})(u)\) （见上面的注 3）。

命题 2.—— 设 A 为维数 d 的局部 Macaulay 环，并设 \(\Omega\) 为一个对偶化 A-模。同态 \(\omega: H_{A}^{d}(A) \to D(\Omega)\) 是双射。

同态 \(\omega\) 是映射的归纳系统 \((\omega_{\mathfrak{a}})_{\mathfrak{a}\in\mathcal{D}_{cs}}\) 的极限，其中

\[
\omega_ {\mathfrak {a}}: \operatorname{Ext} _ {\mathrm{A}} ^ {d} (\mathrm{A} / \mathfrak {a}, \mathrm{A}) \longrightarrow \operatorname{Hom} _ {\mathrm{A}} (\Omega , \operatorname{Ext} _ {\mathrm{A}} ^ {d} (\mathrm{A} / \mathfrak {a}, \Omega))
\]

它把每个元素 \(u\) （属于 \(\operatorname{Ext}_{\mathrm{A}}^{d}(\mathrm{A}/\mathfrak{a},\mathrm{A})\) ）对应到映射 \(x\mapsto f_{x}\circ u\) （A, X, 第 114 页）。因此只需证明每个映射 \(\omega_{\mathfrak{a}}\) 都是双射。

设 \(\mathfrak{a}\) 为 \(\mathcal{D}_{cs}\) 中由对 A 完全割的序列 \(\mathbf{x} = (x_1, \ldots, x_d)\) 所生成的一个理想。Koszul 复形 \(\mathbf{K}^{\bullet}(\mathbf{x}, \mathrm{A})\) 给出 A/ \(\mathfrak{a}\) 的一个投射分解；对每个 A-模 M，A-模 \(\mathrm{H}^d\big(\mathrm{Homgr}_\mathrm{A}(\mathbf{K}^{\bullet}(\mathbf{x}, \mathrm{A}), \mathrm{M})\big)\) 典范地等同于 M/ \(\mathfrak{a}\mathrm{M}\) （A, X, 第 155 页）。由此导出一个同构（A, X, 第 100 页）

\[
\varphi_ {\mathrm{M}}: \mathrm{M} / \mathfrak {a} \mathrm{M} \longrightarrow \operatorname{Ext} ^ {d} (\mathrm{A} / \mathfrak {a}, \mathrm{M}).
\]

设 \(x \in \Omega\) 。注意到同前引第 103 页，命题 2，我们有一个交换图

\[
\begin{array}{c c c} \mathrm{A/a} & \xrightarrow {\varphi_ {\mathrm{A}}} & \operatorname{Ext} _ {\mathrm{A}} ^ {d} (\mathrm{A/a}, \mathrm{A}) \\ \bar {f} _ {x} \Bigg \downarrow & & \Bigg \downarrow \operatorname{Ext} (1 _ {\mathrm{A} / a}, f _ {x}) \\ \Omega / a \Omega & \xrightarrow {\varphi_ {\Omega}} & \operatorname{Ext} _ {\mathrm{A}} ^ {d} (\mathrm{A/a}, \Omega) \end{array}
\]

其中 \(\bar{f}_x\) 是由 \(f_x\) 通过取商导出的同态。由此可知，若对每个 A-模 M，把 \(\mathrm{Ext}_{\mathrm{A}}^{d}(\mathrm{A}/\mathfrak{a},\mathrm{M})\) 与 M/ \(\mathfrak{a}\) M （借助 \(\varphi_{\mathrm{M}}\) ）等同起来，则同态 \(\omega_{\mathfrak{a}}\) 就等同于从 A/ \(\mathfrak{a}\) 到 \(\mathrm{Hom}_{\mathrm{A}}(\Omega,\Omega/\mathfrak{a}\Omega)\) 的、把 1 映为典范满射的那个 A-线性映射，也就是典范映射 A/ \(\mathfrak{a}\longrightarrow\mathrm{End}_{\mathrm{A}/\mathfrak{a}}(\Omega/\mathfrak{a}\Omega)\) 。但由于 A/ \(\mathfrak{a}\) -模 \(\Omega/\mathfrak{a}\Omega\) 是对偶化的（\(\S\) 9, 第 2 节，命题 4），这个映射是双射（ \(\S\) 9, 第 4 节，命题 6），这就证明了本命题。

借助同构 \(\widehat{\alpha}_{\Omega}\) （§ 8, 第 3 节，定理 2, b)），把 Matlis 双对偶 D(D(Ω)) 与 Ω 等同起来。

推论.—— 同态 D(ω): \(\widehat{\Omega} \rightarrow \mathrm{D}(\mathrm{H}_{\mathrm{A}}^{d}(\mathrm{A}))\) 是一个同构。

设 M 为一个 A-模，并设 i 为一个整数。考虑典范同态（§ 8, 第 7 节）

\[
\rho_ {d - i} (\mathrm{M}, \Omega): \mathrm{Tor} _ {d - i} ^ {\mathrm{A}} (\mathrm{M}, \mathrm{D} (\Omega)) \longrightarrow \mathrm{D} \big (\mathrm{Ext} _ {\mathrm{A}} ^ {d - i} (\mathrm{M}, \Omega) \big)
\]

\[
\theta^ {d - i} (\mathrm{M}, \mathrm{H} _ {\mathrm{A}} ^ {d} (\mathrm{A})): \operatorname{Ext} _ {\mathrm{A}} ^ {d - i} \left(\mathrm{M}, \mathrm{D} \left(\mathrm{H} _ {\mathrm{A}} ^ {d} (\mathrm{A})\right)\right) \longrightarrow \mathrm{D} \left(\operatorname{Tor} _ {d - i} ^ {\mathrm{A}} \left(\mathrm{M}, \mathrm{H} _ {\mathrm{A}} ^ {d} (\mathrm{A})\right)\right).
\]

利用同构 \(\omega : \mathrm{H}_{\mathrm{A}}^{d}(\mathrm{A}) \to \mathrm{D}(\Omega)\) 、\(\mathrm{D}(\omega) : \widehat{\Omega} \to \mathrm{D}(\mathrm{H}_{\mathrm{A}}^{d}(\mathrm{A}))\) （命题 2 的推论 1）以及 \(\tau^{i}(\mathrm{M}) : \mathrm{Tor}_{d-i}^{\mathrm{A}}(\mathrm{M}, \mathrm{H}_{\mathrm{A}}^{d}(\mathrm{A})) \longrightarrow \mathrm{H}_{\mathrm{A}}^{i}(\mathrm{M})\) （ \(n^{\circ}\) 2），我们导出 A-模的典范同态

\[
\gamma^ {i} (\mathrm{M}): \mathrm{H} _ {\mathrm{A}} ^ {i} (\mathrm{M}) \longrightarrow \mathrm{D} \left(\mathrm{Ext} _ {\mathrm{A}} ^ {d - i} (\mathrm{M}, \Omega)\right)
\]

\[
\delta^ {i} (\mathrm{M}): \operatorname{Ext} _ {\mathrm{A}} ^ {d - i} (\mathrm{M}, \widehat {\Omega}) \longrightarrow \mathrm{D} \left(\mathrm{H} _ {\mathrm{A}} ^ {i} (\mathrm{M})\right).
\]

定理 1（Grothendieck 对偶）. 设 A 为维数 d 的局部 Macaulay 环，并设 Ω 为一个对偶化 A-模。

\begin{enumerate}
\def\labelenumi{\alph{enumi})}
\item
  A-模 \(\mathrm{H}_{\mathrm{A}}^{d}(\Omega)\) 是一个 Matlis 模；对每个 A-模 P，用 D(P) 表示 Matlis 对偶 \(\mathrm{Hom}_{\mathrm{A}}(\mathrm{P},\mathrm{H}_{\mathrm{A}}^{d}(\Omega))\) 。
\item
  对每个有限型 \(\mathrm{A}\) -模 M 以及任意整数 \(i\) ，典范同态
\end{enumerate}

\[
\gamma^ {i} (\mathrm{M}): \mathrm{H} _ {\mathrm{A}} ^ {i} (\mathrm{M}) \longrightarrow \mathrm{D} \left(\operatorname{Ext} _ {\mathrm{A}} ^ {d - i} (\mathrm{M}, \Omega)\right)
\]

是 Artin A-模之间的一个同构。

\begin{enumerate}
\def\labelenumi{\alph{enumi})}
\setcounter{enumi}{2}
\tightlist
\item
  对每个 A-模 M 和每个整数 i，典范同态
\end{enumerate}

\[
\delta^ {i} (\mathrm{M}): \operatorname{Ext} _ {\mathrm{A}} ^ {d - i} (\mathrm{M}, \widehat {\Omega}) \longrightarrow \mathrm{D} \left(\mathrm{H} _ {\mathrm{A}} ^ {i} (\mathrm{M})\right)
\]

是 \(\widehat{\mathrm{A}}\) -模之间的一个同构。

这由 § 8, 第 7 节，命题 6 得出。

推论.—— 设 A 为一个局部 Macaulay 环，M 为一个维数为 e 的非零有限生成 A-模。A-模 \(\mathrm{H}_{\mathrm{A}}^{e}(\mathrm{M})\) 非零。

根据第 1 节的注 2，我们可以假设局部环 A 是完备的。在此情形 A 具有一个对偶化模 Ω（§ 9, 第 3 节，命题 6 的推论 3）；若 \(\mathrm{H}_{\mathrm{A}}^{e}(M)\) 为零，则它的 Matlis 对偶 \(\operatorname{Ext}_{\mathrm{A}}^{d-e}(M, \Omega)\) 也为零，这与 § 9, 第 1 节，命题 3, b) 相矛盾。

注.—— 1) 当 A-模 M 是有限生成的时，\(\widehat{\mathrm{A}}\) -模 \(\operatorname{Ext}_{\mathrm{A}}^{d-i}(\mathrm{M}, \widehat{\Omega})\) 等同于 \(\widehat{\mathrm{A}} \otimes_{\mathrm{A}} \operatorname{Ext}_{\mathrm{A}}^{d-i}(\mathrm{M}, \Omega)\) （A, X, 第 108 页，命题 7, c)），且 \(\delta^{i}(\mathrm{M})\) 也可由 \(\mathrm{D}(\gamma^{i}(\mathrm{M}))\) 与双对偶性同构复合而得到。

\begin{enumerate}
\def\labelenumi{\arabic{enumi})}
\setcounter{enumi}{1}
\tightlist
\item
  设 \(u: \mathrm{M} \to \mathrm{M}'\) 为 A-模的一个同态。由第 2 节的注 1 以及 § 8, 第 7 节的注，下列诸图是交换的：
\end{enumerate}

\[
\begin{array}{c c c} \mathrm{H} _ {\mathrm{A}} ^ {i} (\mathrm{M}) & \xrightarrow {\gamma^ {i} (\mathrm{M})} & \mathrm{D} (\mathrm{Ext} _ {\mathrm{A}} ^ {d - i} (\mathrm{M}, \Omega)) \\ \mathbb {1} _ {\mathrm{A}} ^ {i} (u) \Bigg \downarrow & & \Bigg \downarrow \mathrm{D} (\mathrm{Ext} (u, 1)) \\ \mathrm{H} _ {\mathrm{A}} ^ {i} (\mathrm{M} ^ {\prime}) & \xrightarrow {\gamma^ {i} (\mathrm{M} ^ {\prime})} & \mathrm{D} (\mathrm{Ext} _ {\mathrm{A}} ^ {d - i} (\mathrm{M} ^ {\prime}, \Omega)) \end{array}
\]

\[
\begin{array}{c c c} \operatorname{Ext} _ {\mathrm{A}} ^ {d - i} (\mathrm{M} ^ {\prime}, \widehat {\Omega}) & \xrightarrow {\delta^ {i} (\mathrm{M} ^ {\prime})} & \mathrm{D} (\mathrm{H} _ {\mathrm{A}} ^ {i} (\mathrm{M} ^ {\prime})) \\ \operatorname{Ext} (u, 1) \Bigg \downarrow & & \Bigg \downarrow \mathrm{D} (\mathrm{H} _ {\mathrm{A}} ^ {i} (u)) \\ \operatorname{Ext} _ {\mathrm{A}} ^ {d - i} (\mathrm{M}, \widehat {\Omega}) & \xrightarrow {\delta^ {i} (\mathrm{M})} & \mathrm{D} (\mathrm{H} _ {\mathrm{A}} ^ {i} (\mathrm{M})) \end{array} .
\]

\begin{enumerate}
\def\labelenumi{\arabic{enumi})}
\setcounter{enumi}{2}
\tightlist
\item
  设
\end{enumerate}

\[
0 \rightarrow \mathrm{M} ^ {\prime} \rightarrow \mathrm{M} \rightarrow \mathrm{M} ^ {\prime \prime} \rightarrow 0\tag{ℰ}
\]

为 A-模的一个正合列。由第 2 节的注 2 以及 § 8, 第 7 节的注，下列诸图是交换的：

\[
\begin{array}{c c c} \mathrm{H} _ {\mathrm{A}} ^ {i - 1} (\mathrm{M} ^ {\prime \prime}) & \xrightarrow {\gamma^ {i - 1} (\mathrm{M} ^ {\prime \prime})} & \mathrm{D} (\mathrm{Ext} _ {\mathrm{A}} ^ {d - i + 1} (\mathrm{M} ^ {\prime \prime}, \Omega)) \\ \partial^ {i - 1} (\mathscr {E}) \Bigg \downarrow & & \Bigg \downarrow (- 1) ^ {d - i + 1} \mathrm{D} (\delta^ {d - i} (\mathscr {E}, \Omega)) \\ \mathrm{H} _ {\mathrm{A}} ^ {i} (\mathrm{M} ^ {\prime}) & \xrightarrow {\gamma^ {i} (\mathrm{M} ^ {\prime})} & \mathrm{D} (\mathrm{Ext} _ {\mathrm{A}} ^ {d - i} (\mathrm{M} ^ {\prime}, \Omega)) \end{array}
\]

\[
\begin{array}{c c c} \operatorname{Ext} _ {\mathrm{A}} ^ {d - i} (\mathrm{M} ^ {\prime}, \widehat {\Omega}) & \xrightarrow {\delta^ {i} (\mathrm{M} ^ {\prime})} & \mathrm{D} (\mathrm{H} _ {\mathrm{A}} ^ {i} (\mathrm{M} ^ {\prime})) \\ \delta^ {d - i} (\mathscr {E}, \widehat {\Omega}) \Bigg \downarrow & & \Bigg \downarrow (- 1) ^ {d - i + 1} \mathrm{D} (\partial^ {i - 1} (\mathscr {E})) \\ \operatorname{Ext} _ {\mathrm{A}} ^ {d - i + 1} (\mathrm{M} ^ {\prime \prime}, \widehat {\Omega}) & \xrightarrow {\delta^ {i - 1} (\mathrm{M} ^ {\prime \prime})} & \mathrm{D} (\mathrm{H} _ {\mathrm{A}} ^ {i - 1} (\mathrm{M} ^ {\prime \prime})) \quad . \end{array}
\]

例.—— 设 A 为维数 1 的 Noether 局部整环；用 K 表示它的分式域。设 Ω 为一个对偶化 A-模，并设 M 为一个有限生成 A-模。A-模 \(\mathrm{H}_{\mathrm{A}}^{0}(M)\) 与 \(\mathrm{H}_{\mathrm{A}}^{1}(M)\) 分别典范地等同于 T(M) 与 \((K/A) \otimes_{A} M\) （第 1 节，例 5）。在这些等同之下，对偶同构

\[
\gamma^ {0} (\mathrm{M}): \mathrm{T} (\mathrm{M}) \longrightarrow \mathrm{D} \left(\operatorname{Ext} _ {\mathrm{A}} ^ {1} (\mathrm{M}, \Omega)\right), \quad \gamma^ {1} (\mathrm{M}): (\mathrm{K} / \mathrm{A}) \otimes_ {\mathrm{A}} \mathrm{M} \longrightarrow \mathrm{D} \left(\operatorname{Hom} _ {\mathrm{A}} (\mathrm{M}, \Omega)\right)
\]

（定理 1）就是 § 9, 第 6 节中定义的那些同构。

